\chapter{Steps, Step Numbers and Results}\label{ch:steps}

A Test Step is the unit of execution and the unit of recorded result. Two commands build one:
\texttt{StepParams(name)} and its setters.

\section{What a step must have}

An ordinary step needs a Step Number and an Action. A Step Name is optional; without one the
Run screen and the report show nothing under the number. The Step Number is free-form text
(\texttt{"1."}, \texttt{"4.0.1"}) shown on the report; it is not the step's position in the
run, and duplicates among siblings are a load error.

These are the setters, and what each one does:

\begin{itemize}
  \item \texttt{:setStepNo(s)} --- the Step Number as displayed.
  \item \texttt{:setEnable(b)} --- the step's default enabled state. The Operator can change
    it before the run; a step disabled at run start never executes and records
    \texttt{Result.NOTRUN}.
  \item \texttt{:setInputs(t)} --- the step's Step Input declarations (chapter~\ref{ch:inputs}).
  \item \texttt{:setAction(fn)} --- the step's logic.
  \item \texttt{:setInclude(file, keyvals)} --- makes the step a Sub-schedule header instead
    (chapter~\ref{ch:sub-schedules}).
  \item \texttt{:setPrefix(s)} --- only with \texttt{setInclude}: prepended to the included
    steps' Step Numbers.
\end{itemize}

\example{../examples/test-writer/steps.tslua}{steps.tslua}

\section{Results}

A step's result is whatever its last \texttt{ATE.report} call said. Every call is logged, only
the last one counts, and a step that never reports records \texttt{Result.NOREPORT} and logs a
WARNING. ATE stores \texttt{actual} and \texttt{expected} as text and does no comparing: the
schedule decides whether a reading passes and says so with the first argument.

These are the seven values a step can record:

\begin{center}
\begin{tabular}{@{}lll@{}}
\toprule
\texttt{Result} & Means & Fails the run \\
\midrule
PASS & the step ran and passed & no \\
FAIL & the step ran and failed & yes \\
ERROR & the step raised, or was aborted, or said it could not be completed & yes \\
SKIPPED & the step ran and chose not to test anything & no \\
NOTRUN & the step never executed & no \\
FORCED\_NOTRUN & the step does not apply to this UUT & no \\
NOREPORT & the step ran but never reported & no \\
\bottomrule
\end{tabular}
\end{center}

A Test Run is \texttt{FAIL} if any step is \texttt{FAIL} or \texttt{ERROR}. \texttt{NOREPORT}
does not fail the run on its own; it is a warning the Operator has to read.

\example{../examples/test-writer/results.tslua}{results.tslua}

Run that schedule and the record shows each of the seven values, with the WARNING that
\texttt{NOREPORT} carries and the run's overall result \texttt{FAIL} because of step 2.

\cando{write a schedule whose steps report every result you need, and explain what each step
left in the record.}
